\documentclass[10pt]{article}
\usepackage[margin=0.75in]{geometry}
\usepackage[T1]{fontenc}
\usepackage{lmodern}
\usepackage{amsmath,amssymb}
\usepackage{enumitem}
\usepackage[colorlinks,linkcolor=blue!50!black,urlcolor=blue!50!black]{hyperref}
\setlist{itemsep=1pt,topsep=2pt}
\setlength{\parindent}{0pt}
\setlength{\parskip}{4pt}
\newcommand{\E}{\mathbb{E}}
\newcommand{\slide}[1]{\subsection*{Slide ``#1''}}
\newcommand{\say}[1]{\par\textit{Say it like this.} #1\par}

\title{Cheat Sheet for the Project Part of Slides v3}
\author{}
\date{6 November 2026}

\begin{document}
\maketitle

This sheet covers the slides from the section ``Research question and contribution'' to ``Questions for discussion''.
For each slide it gives the reading of every symbol, the argument in plain words, a sentence to say and the answers
to likely questions.

\section*{Notation used throughout}

\begin{itemize}
\item $W=(Y,Z)$ is one row of the data, with the outcome $Y$ and $Z=(D,X)$, the binary treatment $D$ and the
covariates $X$. A data set has $n$ independent rows $W_1,\dots,W_n$, all with the law $P$, the data-generating
process.
\item The subscript $0$ marks the true object under $P$. So $g_0$, $\theta_0$, $m_0$ and $\alpha_0$ are the true
outcome regression, the true target, the true propensity score and the true Riesz representer. The $0$ does not mean
that anything is set to zero. The estimator name \texttt{aipw\_cf\_alpha0} means AIPW with the true $\alpha_0$.
\item In $\hat\alpha_{0.3}$ and $\hat\alpha_{0.01}$ the subscript has another meaning. It is the clipping bound $c$
of the estimated propensity score, so $\hat\alpha_{0.3}$ uses $\hat m$ clipped to $[0.3,0.7]$ and $\hat\alpha_{0.01}$
is almost unclipped.
\item $g_0(z)=\E[Y\mid Z=z]$ is the outcome regression, so $g_0(1,x)$ and $g_0(0,x)$ are the mean outcomes of treated
and of control units with covariates $x$. The noise is $\varepsilon=Y-g_0(Z)$, with $\E[\varepsilon\mid Z]=0$.
\item The target is the ATE $\theta_0=\E[g_0(1,X)-g_0(0,X)]$. With $\nu(W;g)=g(1,X)-g(0,X)$ we write
$\theta(g)=\E[\nu(W;g)]$, the ATE that a regression function $g$ implies, so $\theta_0=\theta(g_0)$.
\item $m_0(x)=P(D=1\mid X=x)$ is the propensity score.
\item $\langle a,b\rangle=\E[a(Z)b(Z)]$ and $\|a\|^2=\langle a,a\rangle$. $L^2(P_Z)$ is the set of functions $h$ of
$Z$ with $\E[h(Z)^2]<\infty$, the functions for which these means exist.
\item The Riesz representer $\alpha_0$ is the function with $\theta(g)=\langle\alpha_0,g\rangle$ for every $g$. For
the ATE, $\alpha_0(d,x)=d/m_0(x)-(1-d)/\{1-m_0(x)\}$, the inverse propensity weight with a sign. To check it, note
that $\E[D\,g(D,X)/m_0(X)\mid X]=g(1,X)$, so $\E[\alpha_0(Z)g(Z)]=\E[g(1,X)-g(0,X)]$.
\item $\hat\theta=\hat\theta(Z_{1:n},Y_{1:n})$ is any estimator, for example the trained network, written as a function
of the covariates and treatments $Z_{1:n}=(Z_1,\dots,Z_n)$ and the outcomes $Y_{1:n}=(Y_1,\dots,Y_n)$.
\item $\E[\cdot]$ or $\E_P[\cdot]$ is the mean over many data sets drawn from $P$, so $\E[\hat\theta]-\theta_0$ is the
bias.
\end{itemize}

\section*{Section ``Research question and contribution''}

\slide{Project question and motivation}
A score $\psi(W;\theta,\eta)$ is an estimating function whose mean is zero at the true $\theta_0$ and the true
nuisance $\eta_0$, and Neyman orthogonality says that the mean of the score does not move to first order when the
nuisance moves. A causal PFN returns the estimate in one forward pass and has no score, so this property cannot be
checked for it. In simulations based on clinical data, \textit{Mour\~ao et al.\ (2026)} found that CausalPFN was
biased and that its credible intervals covered poorly. \textit{Melnychuk et al.\ (2026)} prove guarantees for an
idealised network that returns the exact posterior, and the difference between a trained network and this idealised
network, the amortisation error, is outside their theory.
\begin{itemize}
\item Main question in words. When does an estimator without a score have a bias that does not react to small errors in
the outcome regression, and how can we check this with one data set and the network?
\item Sub-question 1 is answered by first-order insensitivity (next section). Sub-question 2 is the diagnostic and the
block test. Sub-questions 3 and 4 concern training targets, priors and network designs, and today only the variants of
the loss touch them.
\end{itemize}
\say{Orthogonality is a property of a score. Our networks have no score, so we need a property of the estimator
itself, and a way to check it from the data.}

\slide{Contribution and positioning}
Read the table row by row, from the left column (what we propose) to the right column (what is new relative to the
closest work).
\begin{enumerate}
\item The criterion is a derivative condition on the bias at a fixed sample size. \textit{Hirano and Porter (2012)}
define local asymptotic unbiasedness for $n\to\infty$, and \textit{Giordano et al.\ (2026)} study local balance of
posterior means.
\item The diagnostic is the imbalance of the implicit weights. Balance checks of weights exist for estimators linear in
the outcomes (\textit{Chattopadhyay and Zubizarreta 2023, Knaus 2024}). New are nonlinear estimators such as networks,
the fact that the expected imbalance equals the sensitivity, and a test with a known level.
\item The weighting representation $\hat\theta=n^{-1}\sum_iw_iY_i$ is known for linear estimators (\textit{Armstrong
and Koles\'ar 2021}), and we show it for nonlinear equivariant estimators.
\item The Bayes rules of in-context learners are theory work that is not on today's slides.
\item The measurement on trained causal PFNs is the experimental part.
\end{enumerate}

\section*{Section ``First-order insensitivity and implicit weights''}

\slide{Outcome shifts and first-order insensitivity}
\textbf{What is $h$?} A function of $(d,x)$, that is, of the treatment and the covariates, with $\E[h(Z)^2]<\infty$. It
gives the shape of a change of the outcome regression. Three examples are useful.
\begin{itemize}
\item $h(d,x)=d$ adds the same amount to every treated outcome. The ATE changes by exactly that amount.
\item $h(d,x)=x_1$ moves the outcomes of both arms in the same way. The ATE does not change, since
$h(1,x)-h(0,x)=0$.
\item $h=\alpha_0$ is the direction in which the ATE changes most per unit size of the change, since
$|\langle\alpha_0,h\rangle|\le\|\alpha_0\|\,\|h\|$ by the Cauchy--Schwarz inequality, with equality for $h$
proportional to $\alpha_0$.
\end{itemize}
\textbf{What is $\delta$?} A real number, the size of the shift. $\delta=0$ is the true law $P$, and we only need small
$\delta$, because the errors of a learned regression are small.

\textbf{What is $P_{\delta h}$?} The law of $(Y+\delta h(Z),Z)$. Take every row, keep $Z$, and add $\delta h(Z)$ to its
outcome. The noise does not change, and the outcome regression becomes $g_0+\delta h$. Since $\theta$ is linear,
\[
\theta(P_{\delta h})=\theta(g_0+\delta h)=\theta_0+\delta\,\theta(h)=\theta_0+\delta\langle\alpha_0,h\rangle,
\]
and for the ATE $\langle\alpha_0,h\rangle=\E[h(1,X)-h(0,X)]$.

\textbf{How to read $b_n(\delta;h)=\E_{\delta h}[\hat\theta]-\theta(P_{\delta h})$.} Read it as ``$b$ n of delta and
h is the expected value of theta hat under P delta h, minus theta of P delta h''. The first term is the average of the
estimate over many data sets of size $n$ drawn from the shifted law, and the second term is the true ATE in that
shifted world. So $b_n(\delta;h)$ is the bias of the estimator in a world whose outcome regression is moved by
$\delta h$. The index $n$ is there because the bias of a learner depends on the sample size, and $b_n(0;h)$ is the
ordinary bias at $P$.

\textbf{What is $\partial_\delta b_n(0;h)$?} The slope of the bias as a function of $\delta$, taken at $\delta=0$, that
is, at the true law. The $0$ inside the bracket is the value of $\delta$ at which the derivative is taken. A slope of
zero means that, to first order, the bias does not change when the regression moves along $h$. Equivalently, the mean
of the estimate moves by $\delta\langle\alpha_0,h\rangle$, as much as the target. A negative slope means that the
estimate moves less than the target, which is what shrinkage does.

\textbf{Relation to Neyman orthogonality.} Orthogonality asks that the mean of the score does not react to first-order
errors in the nuisance. First-order insensitivity asks that the mean of the estimate does not react to first-order
changes of the outcome regression. Both are statements about derivatives at the truth.
\say{We imagine a slightly different world in which the outcome regression is moved by $\delta$ times $h$. The true
ATE then moves by $\delta$ times $\langle\alpha_0,h\rangle$. If the average estimate moves by the same amount, the
estimator is first-order insensitive in direction $h$, and errors in learning the regression along $h$ do not turn into
bias.}

If asked about the last bullet. Local asymptotic unbiasedness looks at shifts of size $\delta=t/\sqrt n$ and asks that
$\sqrt n$ times the bias vanishes as $n\to\infty$; our condition is its derivative part at a fixed $n$.

\slide{Implicit weights}
\textbf{Reading of $w_i=n\,\partial\hat\theta/\partial Y_i$.} If the outcome of row $i$ increases by a small amount $e$,
the estimate increases by about $e\,w_i/n$. The factor $n$ puts the weights on the scale of the weights of an average,
which gives $1/n$ to each row, so they are comparable with inverse propensity weights. Automatic differentiation gives
all $n$ derivatives in one backward pass.

\textbf{Examples.} For inverse propensity weighting, $\hat\theta=n^{-1}\sum_i\hat\alpha(Z_i)Y_i$ and $w_i=\hat\alpha(Z_i)$.
For the difference in means, $w_i=n/n_1$ for treated and $w_i=-n/n_0$ for control rows. For a network the weights
depend on all outcomes, and they are computed at the observed data.

\textbf{Derivation of the formula.} Under the shift, $Y_i$ becomes $Y_i+\delta h(Z_i)$. By the chain rule, the
derivative of $\hat\theta$ in $\delta$ is $\sum_i(\partial\hat\theta/\partial Y_i)h(Z_i)=n^{-1}\sum_iw_ih(Z_i)$. The
mean of the estimate therefore moves by $\E[n^{-1}\sum_iw_ih(Z_i)]$ per unit of $\delta$, the target moves by
$\langle\alpha_0,h\rangle$, and the sensitivity is the difference. The domination condition allows us to exchange the
derivative and the mean.

\textbf{Balance.} The weights balance $h$ when their weighted mean of $h$ equals, in expectation, the change of the
target. The true $\alpha_0$ satisfies $\E[\alpha_0(Z)h(Z)]=\E[h(1,X)-h(0,X)]$ for every $h$. Weights that behave on
average like $\alpha_0$ therefore have zero sensitivity in every direction, and for an estimator that treats the rows
symmetrically this is the condition $\E[w_i\mid Z_i]=\alpha_0(Z_i)$.

\textbf{Plug-in and AIPW.} $\hat Sh$ is the regression method applied to $h$ as if $h$ were the outcome. For the plug-in
estimator the sensitivity is $\E\langle\alpha_0,\hat Sh-h\rangle$, the error the method makes when the truth moves by
$h$, weighted by $\alpha_0$. For AIPW it is $\E\langle\alpha_0-\hat a,\hat Sh-h\rangle$, the product of two errors,
which is why AIPW is less sensitive.

\textbf{Sample imbalance.} $\hat\Delta(h)=n^{-1}\sum_iw_ih(Z_i)-n^{-1}\sum_i\nu(W_i;h)$ is the weighted mean of $h$ minus
the sample version of the change of the ATE, $n^{-1}\sum_i\{h(1,X_i)-h(0,X_i)\}$. Its mean is the sensitivity, and it
needs only the weights and $h$, neither $g_0$ nor $\alpha_0$.

\textbf{Why $S_\alpha=\E[\hat\Delta(\alpha_0)]/\|\alpha_0\|^2$.} Along $h=\alpha_0$ the target moves by
$\delta\langle\alpha_0,\alpha_0\rangle=\delta\|\alpha_0\|^2$, and the mean of the estimate moves by
$\delta(\|\alpha_0\|^2+\E[\hat\Delta(\alpha_0)])=\delta\|\alpha_0\|^2(1+S_\alpha)$. So $1+S_\alpha$ is the share of the
change of the target that the estimate follows. $S_\alpha=0$ means full response, $S_\alpha=-0.2$ means 80\%, and
$S_\alpha=-1$ means no response at all. It needs the true $\alpha_0$, so we compute it only in simulations.
\say{The weights tell us how the estimate reacts to each outcome. If they act on average like inverse propensity
weights, the estimate follows every shift of the regression, and if not, the gap is the sensitivity.}

\slide{Estimators that are equivariant in the outcomes}
\textbf{Why it matters.} Weights are derivatives, so in general they only describe small changes. For equivariant
estimators the estimate is exactly the weighted mean $n^{-1}\sum_iw_iY_i$, so the weights describe the whole estimate
and its bias.

\textbf{Reading of the definition.} Add $a$ to every outcome and the estimate moves by $a\kappa_n$; multiply every
outcome by $b>0$ and the estimate is multiplied by $b$. For the ATE, adding a constant to all outcomes changes no
treatment effect, so $\kappa_n=0$.

\textbf{Why our networks are equivariant.} They compute $\tilde Y_i=(Y_i-\bar Y)/s_Y$ from the context, predict on this
scale and multiply the result by $s_Y$. Adding $a$ leaves $\tilde Y$ unchanged, so the estimate does not change, and
multiplying by $b$ leaves $\tilde Y$ unchanged and multiplies $s_Y$ by $b$, so the estimate is multiplied by $b$.

\textbf{Euler's theorem.} From $\hat\theta(bY)=b\,\hat\theta(Y)$, the derivative in $b$ at $b=1$ gives
$\sum_iY_i\,\partial\hat\theta/\partial Y_i=\hat\theta$, which is $\hat\theta=n^{-1}\sum_iw_iY_i$.

\textbf{The decomposition.} Insert $Y_i=g_0(Z_i)+\varepsilon_i$, and add and subtract $n^{-1}\sum_i\nu(W_i;g_0)$.
\begin{itemize}
\item Imbalance along $g_0$. $\hat\Delta(g_0)=n^{-1}\sum_iw_ig_0(Z_i)-n^{-1}\sum_i\nu(W_i;g_0)$. Its mean is the
sensitivity in the direction $h=g_0$, the direction that scales the regression and the target together.
\item Mean zero. $n^{-1}\sum_i\nu(W_i;g_0)-\theta_0$ is sampling noise, the same for every estimator, because
$\E[\nu(W;g_0)]=\theta_0$.
\item Noise term. $n^{-1}\sum_iw_i\varepsilon_i$. When the weights depend only on $Z_{1:n}$, its mean is zero because
$\E[\varepsilon_i\mid Z_{1:n}]=0$. For a network the weights depend on the outcomes and therefore on the noise, and
the mean is the covariance of weights and noise. For normal noise it equals the curvature
$\sum_i\E[\sigma^2(Z_i)\partial^2\hat\theta/\partial Y_i^2]$, which is zero for linear estimators.
\end{itemize}
Taking means gives bias $=$ sensitivity along $g_0$ $+$ mean of the noise term. For the small networks the bias is
almost all imbalance along $g_0$ (correlation 0.98 to 0.99), so their noise term is small.
\say{Because the networks standardise the outcomes, their estimate is exactly a weighted mean of the outcomes. Then the
error splits into an imbalance part, which our diagnostic targets, a pure sampling part and a curvature part.}

\slide{A diagnostic from one data set}
\textbf{Goal.} Test the hypothesis that the sensitivity along a direction $\hat h$ is zero, with one data set and the
network, and without $g_0$ or $\alpha_0$.

\textbf{The imbalance for the ATE.} The first sum is the weighted mean of $\hat h$ with the weights of the network. The
second sum is how much the ATE would change if the regression moved by $\hat h$. Under the hypothesis they agree in
expectation.

\textbf{Why a direction close to $\alpha_0$.} The sensitivity $S_\alpha$ comes mostly from rows with weak overlap, where
$|\alpha_0|$ is large and the networks shrink most. A test along an unrelated direction, such as $h=x_1$, would have
little power.

\textbf{The split.} The last $n_s=\max(50,\lceil 0.2n\rceil)$ rows, which are 50 rows at $n=250$ and 100 rows at
$n=500$, are used only to fit the direction. The network is run on the other $n_e$ rows (200 or 400). Given the split,
$\hat h$ is a fixed function for the evaluation rows, so the mean of the imbalance is the sensitivity along $\hat h$.
If $\hat h$ were fitted on the same rows, the data would be used twice and this would fail.

\textbf{The clipping.} $\hat m$ is clipped to $[0.3,0.7]$, so $|\hat\alpha_{0.3}|\le1/0.3\approx3.3$. Without it, rows
with $\hat m$ near 0 or 1 get huge weights and dominate the imbalance.

\textbf{Choice of $\hat m$.} Among logistic regression, gradient boosting and the propensity head of the network, we take
the one whose representer $a$ has the smallest Riesz loss $n_s^{-1}\sum\{a(Z_i)^2-2[a(1,X_i)-a(0,X_i)]\}$ on the
split. The mean of this loss is $\|a-\alpha_0\|^2-\|\alpha_0\|^2$, so it is smallest for the candidate closest to
$\alpha_0$.

\slide{Block test of zero sensitivity}
\textbf{Why blocks.} One run gives one imbalance, and there is no formula for its variance for a network. We cut the
$n_e$ evaluation rows into $q$ disjoint blocks, for example 5 blocks of 80 rows at $n=500$, run the network on each
block alone as its own context, and compute one imbalance per block. Because the rows are independent and the blocks
are disjoint, the $q$ imbalances are independent given the split, so their spread estimates the uncertainty.
Each block imbalance has as its mean the sensitivity of the network at the block size, $n_e/q$ rows.

\textbf{The test.} $t_q=\sqrt q\,\bar\Delta/s_\Delta$ is the usual one-sample $t$-statistic of the $q$ block values. We
reject at 5\% when $|t_q|$ exceeds $t_{q-1,0.975}$, which is 3.18 for $q=4$ and 2.78 for $q=5$.
\textit{Ibragimov and M\"uller (2010)} show that with normal block values this test keeps its level for $2\le q\le14$,
also when the blocks have different variances.

\textbf{Skewness and Hall.} Each row contributes $w_i\hat h(Z_i)-\nu(W_i;\hat h)$, and these contributions contain
inverse propensities, so they are skewed to the right. With few rows per block, the mean and the standard deviation
of the blocks are positively correlated, the $t$-statistic is skewed to the left, and the test rejects too often in the
lower tail. \textit{Hall (1992)} gives a transformation of $t_q$ that uses the estimated skewness $\gamma$ and removes
this skewness to first order, $t_q+\gamma(t_q^2/3+1/6)/\sqrt q+\gamma^2t_q^3/(27q)$.

\textbf{How we check level and power in simulations.} The level is the rejection rate when the hypothesis is true. We
use two estimators whose sensitivity is zero by construction.
\begin{itemize}
\item AIPW with the true $\alpha_0$ and the outcome regression $\hat g$ of the network, cross-fitted. Its mean is
$\theta(\hat g)+\langle\alpha_0,g-\hat g\rangle=\langle\alpha_0,g\rangle$ for every true regression $g$, so it is
unbiased in every shifted world, its bias is always zero, and so is its slope. This is why we use the true $\alpha_0$.
\item AIPW with both true nuisances. Its weights are exactly $\alpha_0(Z_i)$, so it balances every direction.
\end{itemize}
Both use the network's data and rows, so the test sees realistic contributions while the hypothesis holds. The power is
the rejection rate for the plug-in estimate of the network, which is sensitive.
\say{We split the data into a few independent blocks, compute the imbalance in each block and run a $t$-test across the
blocks. Hall's correction removes the skewness that inverse propensities create.}

\section*{Section ``Experiments with trained networks''}

\slide{Networks, priors and truths}
A causal PFN is trained on many simulated tasks drawn from a prior. Each task is a data-generating process, and the
network learns to predict the CATE of the rows from a context of rows. Our plug-in estimate of the ATE is the mean of
the predicted CATEs over the sample. The twelve small networks are 4 variants of the loss times 3 seeds (V0 no
nuisance heads, V1 heads for $m$, $g$ and $\alpha$ supervised with the exact nuisances, V2 additionally an AIPW layer
with cross-fitting in context, V3 head losses annealed to zero by the middle of training).
\begin{itemize}
\item $\rho$ is the correlation between the baseline outcome and the logit of the propensity score, which creates
confounding. The prior with strong confounding uses $|\rho|\in[0.7,0.9]$ and a propensity logit scaled by 1.5 to 2.5,
which gives weak overlap. Both ranges lie inside the training prior, so these truths are a hard region of the prior
and not new territory.
\item A truth is one fixed data-generating process. We have 16 truths per prior with 32 data sets at each $n$, and for
the diagnostic 40 (128 for the AIPW estimate with true nuisances) new truths per prior drawn with another seed.
\item Units are standard deviations of $Y$, so a bias of 0.2 is 0.2 standard deviations.
\item Released networks are black boxes, so their sensitivity is the central finite difference
$\{\hat\theta(Y+\eta h)-\hat\theta(Y-\eta h)\}/(2\eta)$ of the estimate.
\end{itemize}

\slide{Sensitivity of trained networks along the Riesz representer}
Top panels show $S_\alpha$, bottom panels the RMS bias, the root mean square over truths of the bias at each truth,
against $n$. At $n=500$, $S_\alpha=-0.21$ on the training prior and $-0.39$ under strong confounding, so the average
estimate follows only 79\% and 61\% of a shift along $\alpha_0$. The values hardly change from $n=100$ to 500, so the
shrinkage does not disappear with more rows in this range. Strong confounding is worse because the networks shrink
most where the overlap is weak.
\say{The small networks follow only about 80\% of a change of the true effect on the training prior, and about 60\%
under strong confounding.}

\slide{Released causal PFNs}
Under independent mechanisms the prior draws the treatment mechanism independently of the outcome mechanism. The
treatments then carry no information about the outcome regression, and the posterior mean of the ATE is the ATE of the
posterior mean regression $\hat g_\Pi$. Its error contains $\langle\alpha_0,\hat g_\Pi-g_0\rangle$, the plug-in error
that AIPW would correct. CausalPFN follows 83\% and 70\% of a shift along $\alpha_0$, with RMS bias 0.11 and 0.22, and
CausalFM and Do-PFN follow only 26\% to 65\%, with RMS bias 0.46 to 0.83. Our small networks, with $S_\alpha$ of
$-0.21$ and $-0.39$, are slightly more sensitive than CausalPFN.

\slide{Noise term and corrections}
$\hat\theta^{\mathrm c}=\hat\theta-\hat\Delta(\hat g)$ is a plug-in with $\hat g$ plus the residuals $Y_i-\hat g(Z_i)$
weighted by the network's own weights, so it looks like AIPW with the network's weights in place of $\hat\alpha$. Its
bias is $\langle\omega_n-\alpha_0,g_0-\hat g\rangle$ plus the noise term, where $\omega_n(z)=\E[w_1\mid Z_1=z]$ is the
average weight. This is a product of a weight error and a regression error. With $\hat g$ from gradient boosting fitted
on other rows, the RMS bias of the small networks at $n=500$ falls from 0.215 to 0.091 and from 0.344 to 0.164.
\say{The network's own weights can be used to correct its estimate, and the correction halves the bias of the small
networks.}

\slide{Level of the block test}
Each bar is the larger of the two rejection rates at $n=250$ and $500$ for the two estimators with zero sensitivity,
pooled over both priors, along four directions. The dotted line is the nominal 5\%. Along the almost unclipped
representer $\hat\alpha_{0.01}$ and with the bound 0.1, the test rejects too often, in the lower tail, because the
contributions are strongly skewed (median skewness 2.9 to 7.7). With the bound 0.3, chosen on 24 separate truths per
prior, the rates are 4.7\% to 6.4\% pooled and at most 6.8\% on one prior. The direction $h_{\mathrm o}=(2d-1)\hat
m(1-\hat m)$ is the overlap direction, which is bounded and puts weight on rows with good overlap.

\slide{Estimand of the diagnostic}
The share along a direction is the imbalance divided by the change of the target along that direction, the analogue of
$S_\alpha$ for that direction. The bars are correlations over truths between its mean and $S_\alpha$. Along the true
$\alpha_0$ the share is $S_\alpha$ itself up to sampling noise, so the correlation is at least 0.99. Along the fitted
directions it is at most 0.55 and 0.24, because $S_\alpha$ comes mostly from rows with weak overlap, which the clipping
downweights. So the diagnostic does not estimate $S_\alpha$, and we use it only as a test of zero sensitivity.

\slide{Results of the diagnostic on the small networks}
The middle column is the block test with $q=4$ without the transformation on the 16 truths per prior. The right column
is the block test with $q=5$ after Hall's transformation on new truths that played no role in choosing the test, which
guards against a test that looks good by chance.
\begin{itemize}
\item Level. 3.9\% to 5.2\% pooled and at most 5.3\% on one prior, close to the nominal 5\%.
\item Power. The test detects the sensitivity of the plug-in estimate in 58\% to 70\% of the data sets with at most 500
rows. It is lower at $n=500$ than at $n=250$, possibly because the sensitivity per block is smaller with larger blocks,
and this is open.
\item Without the transformation, the $q=4$ test rejects in 7.8\% pooled and 9.0\% under strong confounding on the new
truths, so the correction is needed.
\item Finite differences agree with the backward pass in 99\% of 384 pairs, so the test also works for black-box
networks.
\end{itemize}

\slide{Status and next steps}
The theory (definitions, propositions and proofs) is drafted. The level result for the block test holds for normal block
values and the untransformed statistic, and the level after the transformation rests on the simulations. Next come
larger networks, the power up to $n=2{,}000$ and the corrections. The paper targets ICML 2027, with the ATE in the main
text. IHDP is a semi-synthetic benchmark with known outcome surfaces, and the RealCause version of the LaLonde PSID data
is a simulation fitted to real data, so in both the true effect is known.

\slide{Questions for discussion}
\begin{enumerate}
\item The shift $h(d,x)=d$ asks whether the estimate follows when every treated outcome is one unit higher, which is
easy to explain to users. Shifts of the effect in one covariate are another candidate.
\item The test plays the role of a balance check after matching, run before the estimate is reported.
\item Hall's transformation works in our simulations. A bootstrap-$t$ is the alternative, and theory for both is open.
\item IHDP and LaLonde are planned. ACIC data are a further candidate.
\end{enumerate}

\section*{Likely questions and short answers}
\begin{itemize}
\item \textit{Why not just compute the bias?} It needs the true effect, which is unknown in an application. The
sensitivity can be tested from the data.
\item \textit{Is an insensitive estimator unbiased?} Insensitivity is about the slope of the bias. For an equivariant
estimator, insensitivity along $g_0$ makes it unbiased up to the noise term.
\item \textit{Does running the network on blocks change the estimator?} Yes, each block is a smaller context. The test
concerns the sensitivity at the block size, which is the honest statement.
\item \textit{Why not the bootstrap?} A bootstrap needs many forward passes and theory for a nonlinear estimator, while
independent blocks give a $t$-test that keeps its level with few groups.
\item \textit{Why clip at 0.3?} Clipping trades power for level. The bound was chosen on separate truths.
\item \textit{Why the true $\alpha_0$ in the null estimators?} AIPW with the true $\alpha_0$ is unbiased for every
outcome regression, so its sensitivity is zero by construction, and its rejection rate is the level.
\item \textit{Are the networks out of distribution under strong confounding?} No, both parameters lie inside the
ranges of the training prior.
\end{itemize}

\end{document}
